\section{Avenues for further research}
\label{sec:future}

Each avenue below addresses the limitation carrying the same number in
Section~\ref{sec:limitations}.

\begin{description}[leftmargin=2.4em,style=nextline]

  \item[FR1] Researchers should test whether the accounts carry information at a level of
  aggregation where loan composition does not already dominate, such as the county or the
  program, rather than at the level of an individual intermediary where we have shown they do
  not. The two nulls we report both bite at bank level, and the reason they bite there, that the
  measure is largely a restatement of loan mix across institutions, does not obviously apply to a
  geography whose employers differ while its loan mix does not.

  \item[FR2] Researchers should build the off-balance-sheet side of the AI claim stock from
  securitization prospectuses, vehicle-level filings and supervisory data on lending secured
  against computing hardware, so that side can be measured rather than bounded, and the holder
  comparison can be made on two measured stocks rather than one measured and one bounded.

  \item[FR3] Researchers should separate realized capital gains from wage income inside the
  individual income tax base and recompute the Treasury coefficient on the separated series,
  which would replace the band we carry with a point and tighten every fiscal magnitude that
  depends on it.

  \item[FR4] Researchers should identify the mortgage holder remainder from loan-level sources
  rather than leaving it as an arithmetic residual treated as private in full, which would
  establish whether the federal share we report is conservative by a little or by a lot.

  \item[FR5] Researchers should reconstruct income by source for the historical years from tax
  and survey microdata, which would let the whole series be placed on the adopted wage basis and
  remove the assumption that the two constructions stay close going backwards.

  \item[FR6] Researchers should define the universe and the normalization for the bottom-quintile
  cell explicitly and recompute it, so the one withdrawn distributional magnitude can be restored
  or ruled out rather than left absent.

  \item[FR7] Researchers should extend the novelty search into the archives this one did not
  reach. The protocol is published and the search has been run against a general web index with
  a stated kill criterion, which is what converts ``none located'' from an admission into a
  result. What it does not reach is paywalled journal archives searched systematically rather
  than by query, non-English literature, statistical agency working papers outside the web index
  and unpublished central bank work. A librarian-assisted search across those four would settle
  the question in a way a web search cannot.

  \item[FR8] Researchers should construct labor backing accounts for a second country on the same
  basis. The United States is one observation, and whether a federal share near four fifths is a
  property of this state or of states with large mortgage guarantee programs generally cannot
  be settled from one country. Our own attempt at a euro-area counterpart did not reach the
  standard of the US series, because the published counterpart detail does not support the
  issuer-side leg, so this needs national sources country by country rather than a single
  harmonized release.

  \item[FR9] Researchers should put the whole accounts on a single vintage of the federal
  receipts aggregates. The Treasury coefficient is currently computed on the 2025 national
  accounts while the capital-gains band and the quoted receipts base come from a different
  vintage, so the two differ by about two percentage points for reasons that have nothing to do
  with the construction. One vintage throughout would remove that.

\end{description}
